% GENERATED FILE. Edit canonical lessons, then run npm run generate:books.
% canonical-source-hash: fnv1a64:41c66612d093a35f
% canonical-lessons: RU-C245-strannyy, RU-C245-smeshnoy, RU-C245-prekrasnyy, RU-C245-uzhasnyy, RU-C245-otlichnyy

\chapter{Strange, Funny, Wonderful, Terrible, Excellent}
\label{ch:ru-strange-funny-wonderful-terrible-excellent}
\hlchaptermodality{245}

\begin{chapteropening}
\textbf{By the end of this chapter:} \emph{I can say strange, funny, wonderful, terrible and excellent.}

Complete the last lesson of chapter 245: I can say strange, funny, wonderful, terrible and excellent.
\end{chapteropening}

\section[\texorpdfstring{\ru{странный} (stránnyy)}{странный (stránnyy)}]{\ru{странный} (stránnyy) — strange}

\label{lesson:RU-C245-strannyy}



\begin{quote}
Before the new one: say the Russian for dark, then the Russian for light.
\end{quote}

\subsection*{You'll want to know: \ru{странный}}
\textbf{\ru{странный}} (\emph{stránnyy}) — \textquotedblleft{}strange\textquotedblright{}.

\begin{cousinweb}
\textbf{\ru{странный}} — Old Church Slavonic \emph{stranĭnŭ}, \textquotedblleft{}foreign\textquotedblright{}, from \emph{\ru{страна}}, \textquotedblleft{}a land, a side\textquotedblright{} (Vasmer): strange was first \textquotedblleft{}from another land\textquotedblright{}.
\end{cousinweb}

\begin{grammarlens}[title={What you've built}]
One more word for this chapter.
\end{grammarlens}

\subsection*{Your turn}
\begin{itemize}
\raggedright
  \item \emph{Say it:} \emph{stránnyy}
  \item \emph{Say it:} \emph{stránnyy}, once more
  \item \emph{Say it:} say \emph{stránnyy}
  \item \emph{Recall:} say the Russian for dark, then the Russian for light, then say \emph{stránnyy} again
\end{itemize}

\begin{tcolorbox}[breakable,colback=teal!4,colframe=teal!35!black,title={Before you move on}]
What does \ru{странный} mean? (\textquotedblleft{}Strange\textquotedblright{}.) Say it once more.
\end{tcolorbox}
\section[\texorpdfstring{\ru{смешной} (smeshnóy)}{смешной (smeshnóy)}]{\ru{смешной} (smeshnóy) — funny}

\label{lesson:RU-C245-smeshnoy}



\begin{quote}
Before the new one: say the Russian for light, then the Russian for strange.
\end{quote}

\subsection*{You'll want to know: \ru{смешной}}
\textbf{\ru{смешной}} (\emph{smeshnóy}) — \textquotedblleft{}funny\textquotedblright{}.

\begin{cousinweb}
\textbf{\ru{смешной}} comes from \emph{\ru{смех}}, \textquotedblleft{}laughter\textquotedblright{}, with \emph{-\ru{ной}}: what makes people laugh.
\end{cousinweb}

\begin{grammarlens}[title={What you've built}]
One more word for this chapter.
\end{grammarlens}

\subsection*{Your turn}
\begin{itemize}
\raggedright
  \item \emph{Say it:} \emph{smeshnóy}
  \item \emph{Say it:} \emph{smeshnóy}, once more
  \item \emph{Say it:} say \emph{smeshnóy}
  \item \emph{Recall:} say the Russian for light, then the Russian for strange, then say \emph{smeshnóy} again
\end{itemize}

\begin{tcolorbox}[breakable,colback=teal!4,colframe=teal!35!black,title={Before you move on}]
What does \ru{смешной} mean? (\textquotedblleft{}Funny\textquotedblright{}.) Say it once more.
\end{tcolorbox}
\section[\texorpdfstring{\ru{прекрасный} (prekrásnyy)}{прекрасный (prekrásnyy)}]{\ru{прекрасный} (prekrásnyy) — wonderful}

\label{lesson:RU-C245-prekrasnyy}



\begin{quote}
Before the new one: say the Russian for strange, then the Russian for funny.
\end{quote}

\subsection*{You'll want to know: \ru{прекрасный}}
\textbf{\ru{прекрасный}} (\emph{prekrásnyy}) — \textquotedblleft{}wonderful\textquotedblright{}.

\begin{cousinweb}
\textbf{\ru{прекрасный}} is \emph{\ru{пре}-}, \textquotedblleft{}very\textquotedblright{}, with \emph{\ru{красный}}, which first meant \textquotedblleft{}beautiful\textquotedblright{}, from \emph{\ru{краса}}, \textquotedblleft{}beauty\textquotedblright{} (Vasmer).
\end{cousinweb}

\begin{grammarlens}[title={What you've built}]
One more word for this chapter.
\end{grammarlens}

\subsection*{Your turn}
\begin{itemize}
\raggedright
  \item \emph{Say it:} \emph{prekrásnyy}
  \item \emph{Say it:} \emph{prekrásnyy}, once more
  \item \emph{Say it:} say \emph{prekrásnyy}
  \item \emph{Recall:} say the Russian for strange, then the Russian for funny, then say \emph{prekrásnyy} again
\end{itemize}

\begin{tcolorbox}[breakable,colback=teal!4,colframe=teal!35!black,title={Before you move on}]
What does \ru{прекрасный} mean? (\textquotedblleft{}Wonderful\textquotedblright{}.) Say it once more.
\end{tcolorbox}
\section[\texorpdfstring{\ru{ужасный} (uzhásnyy)}{ужасный (uzhásnyy)}]{\ru{ужасный} (uzhásnyy) — terrible}

\label{lesson:RU-C245-uzhasnyy}



\begin{quote}
Before the new one: say the Russian for funny, then the Russian for wonderful.
\end{quote}

\subsection*{You'll want to know: \ru{ужасный}}
\textbf{\ru{ужасный}} (\emph{uzhásnyy}) — \textquotedblleft{}terrible\textquotedblright{}.

\begin{cousinweb}
\textbf{\ru{ужасный}} is \emph{\ru{ужас}}, \textquotedblleft{}horror\textquotedblright{}, with \emph{-\ru{ный}}. Vasmer calls \emph{\ru{ужас}} \textquotedblleft{}a difficult word\textquotedblright{} and thinks it possibly related to \emph{\ru{гасить}}, \textquotedblleft{}to put out\textquotedblright{}, and Greek \emph{sbénnymi}, \textquotedblleft{}I quench\textquotedblright{}: the terror that chokes.
\end{cousinweb}

\begin{grammarlens}[title={What you've built}]
One more word for this chapter.
\end{grammarlens}

\subsection*{Your turn}
\begin{itemize}
\raggedright
  \item \emph{Say it:} \emph{uzhásnyy}
  \item \emph{Say it:} \emph{uzhásnyy}, once more
  \item \emph{Say it:} say \emph{uzhásnyy}
  \item \emph{Recall:} say the Russian for funny, then the Russian for wonderful, then say \emph{uzhásnyy} again
\end{itemize}

\begin{tcolorbox}[breakable,colback=teal!4,colframe=teal!35!black,title={Before you move on}]
What does \ru{ужасный} mean? (\textquotedblleft{}Terrible\textquotedblright{}.) Say it once more.
\end{tcolorbox}
\section[\texorpdfstring{\ru{отличный} (otlíchnyy)}{отличный (otlíchnyy)}]{\ru{отличный} (otlíchnyy) — excellent}

\label{lesson:RU-C245-otlichnyy}



\begin{quote}
Before the new one: say the Russian for wonderful, then the Russian for terrible.
\end{quote}

\subsection*{You'll want to know: \ru{отличный}}
\textbf{\ru{отличный}} (\emph{otlíchnyy}) — \textquotedblleft{}excellent\textquotedblright{}.

\begin{cousinweb}
\textbf{\ru{отличный}} comes from \emph{\ru{отличить}}, \textquotedblleft{}to set apart\textquotedblright{}, on the root of \emph{\ru{лик}}, \textquotedblleft{}a face, an image\textquotedblright{}: what stands out from the rest.
\end{cousinweb}

\begin{grammarlens}[title={What you've built}]
One more word for this chapter.
\end{grammarlens}

\subsection*{Your turn}
\begin{itemize}
\raggedright
  \item \emph{Say it:} \emph{otlíchnyy}
  \item \emph{Say it:} \emph{otlíchnyy}, once more
  \item \emph{Say it:} say \emph{otlíchnyy}
  \item \emph{Recall:} say the Russian for wonderful, then the Russian for terrible, then say \emph{otlíchnyy} again
\end{itemize}

\begin{tcolorbox}[breakable,colback=teal!4,colframe=teal!35!black,title={Before you move on}]
What does \ru{отличный} mean? (\textquotedblleft{}Excellent\textquotedblright{}.) Say it once more.
\end{tcolorbox}
